\SourceLeaf{285}%
\(GF\) in \(G\) bedeutet); imgleichen die (immer negativen) Werthe von \(\tfrac{206265''\mathrm{d}m}{2m\mathrm{d}S}\) in denselben Punkten mit \(-k^0,-k'\), so wird
\[
\begin{aligned}
206265''l^0&=-2k^0\sin V^0\\
206265''l'&=+2k'\sin V'
\end{aligned}
\]

Die im vorhergehenden Artikel gegebenen Ausdrücke für \(\psi^0,\psi'\), in Secunden verwandelt, werden daher, wenn man die von der Einheit hier nur unmerklich abweichenden Factoren \(\BareRoot[6]{\tfrac{m^0}{m'}},\BareRoot[6]{\tfrac{m'}{m^0}}\) weglässt,
\[
\begin{aligned}
\psi^0&=-\tfrac{1}{3}h(2k^0\sin V^0-k'\sin V')\\
\psi'&=-\tfrac{1}{3}h(2k'\sin V'-k^0\sin V^0)
\end{aligned}
\]

Die dieser Abhandlung beigefügte Tafel gibt in der letzten Columne unter der Ueberschrift \(k\) die Werthe von \(k^0,k'\) für die entsprechenden Werthe von \(S\), die in der ersten Columne unter der Ueberschrift \(Q+q\) aufzusuchen sind; da \(k\) immer positiv ist, und \(\sin V^0,\sin V'\) immer entgegengesetzte Zeichen haben, so wird \(\psi^0\) negativ, \(\psi'\) positiv, wenn \(G\) westlich von \(F\) liegt und umgekehrt: bei der Berechnung erinnere man sich, dass in diesen Formeln \(h\) als in Theilen des Halbmessers ausgedrückt verstanden wird, also der in irgend einem Längenmaasse gegebene Abstand der Punkte \(F,G\) zuvor mit dem in gleichem Maasse ausgedrücktem Werthe von \(A\) zu dividiren ist.

\begingroup\rightskip=0pt plus 1em
Da in unserer conformen Übertragung der Ellipsoidfläche auf die Kugelfläche ein Meridian auf jener wiederum durch einen Meridian auf dieser dargestellt wird, so ist klar, dass jedes Element von \(L\) dieselbe Neigung gegen den Meridian hat wie das entsprechende Element von \(M\), und dass folglich die Azimuthe der geodaetischen Linie in ihren beiden Endpunkten resp. \(V^0+\psi^0\) und \(V'+\psi'\) sein werden: sind aber umgekehrt diese gegeben, so werden sie auf die Kugelfläche reducirt durch Anbringung von \(-\psi^0,-\psi'\). und für die Berechnung dieser stets fast ganz verschwindenden Reductionen ist es offenbar ganz gleichgültig, wenn man in den obigen Formeln anstatt \(V^0,V'\) die Azimuthe auf dem Ellipsoid anwendet.\par\endgroup
\Article{14}
Um nach den gegebenen Vorschriften die Reductionen der Richtungen, behuf der Übertragung vom Ellipsoid auf die Kugel oder umgekehrt, berechnen zu
